\section{Problem Formulation}

\subsection{Vehicle-Cloud Collaborative Anti-Rollover Workflow}

The considered system consists of one autonomous tank truck, surrounding human-driven or automated vehicles, roadside sensing units, and a cloud-side predictive planning service. The target operating domain is a low-penetration highway environment in which the cloud can provide wide-area predictive guidance, but the ego vehicle must retain an independent safety loop. The hierarchy is organized as three nested safety layers. The cloud-side outer loop reduces induced risks by anticipating long-horizon traffic interactions. The vehicle-side planning loop provides degraded-safety fallback when cloud guidance is unavailable or locally unsafe. The high-frequency control loop suppresses state risks caused by roll motion, liquid sloshing, and actuator-limited maneuvers. The ego vehicle is assumed to have onboard perception, localization, state estimation, trajectory replanning, and actuator-level control capability. The cloud-side service is assumed to have access to road topology, fused surrounding-vehicle states, and the ego state through vehicle-road-cloud communication.

At discrete cloud-planning time $k_c$, the cloud constructs a directed traffic graph
\begin{equation}
  \mathcal{G}_{k_c}=(\mathcal{V}_{k_c},\mathcal{E}),
\end{equation}
where $\mathcal{V}_{k_c}$ contains road-node static attributes, surrounding-vehicle dynamic attributes, ego occupancy, and risk-field features, and $\mathcal{E}$ contains road-topology and lane-connectivity attributes. The cloud-side input is
\begin{equation}
  \mathcal{I}_c=\{\mathcal{G}_{k_c},\mathcal{M},x_e(k_c),\mathcal{R}_{nav}\},
\end{equation}
where $\mathcal{M}$ is the road topology, $x_e$ is the ego state, and $\mathcal{R}_{nav}$ is the route or transportation-task constraint. The cloud-side output is
\begin{equation}
  \mathcal{O}_c=\{d_c,\tau_c\},
\end{equation}
where $d_c$ is a long-horizon driving decision, including acceleration and lane-change intention, and $\tau_c=\{r_i^c\}_{i=1}^{N_c}$ is a rollover-aware reference trajectory.

The vehicle-side planning layer receives $\tau_c$, local perception $\mathcal{P}_v$, and the current ego state. It either accepts the cloud trajectory or generates a local emergency trajectory
\begin{equation}
  \tau_v=\{r_i^v\}_{i=1}^{N_v}
\end{equation}
when any of the following events occurs: a sudden obstacle is detected, the cloud trajectory becomes locally infeasible, the predicted rollover index approaches the safety boundary, or the cloud message is delayed, lost, or inconsistent with local perception. The high-frequency control layer tracks the selected trajectory $\tau_r\in\{\tau_c,\tau_v\}$ and outputs steering, braking, and speed commands
\begin{equation}
  u=[\delta,\ p_{brk}^{L1},\ p_{brk}^{R1},\ p_{brk}^{L2},\ p_{brk}^{R2},\ v_{des}]^T.
\end{equation}

\begin{table}[t]
  \centering
  \caption{Layered Information Flow and Time Scale}
  \label{tab:workflow_timescale}
  \begin{tabularx}{\linewidth}{p{0.21\linewidth}p{0.21\linewidth}p{0.26\linewidth}X}
    \toprule
    Layer & Typical rate & Main input & Main output \\
    \midrule
    Cloud-side predictive planning & 1--2 Hz & Traffic graph, road topology, ego state & Long-horizon decision and rollover-aware reference \\
    Vehicle-side emergency planning & about 10 Hz & Cloud reference, local perception, ego state & Locally feasible emergency trajectory \\
    High-frequency anti-rollover control & 20 Hz or higher & Reference trajectory, vehicle-liquid states & Steering, braking, and speed commands \\
    \bottomrule
  \end{tabularx}
\end{table}

Table~\ref{tab:workflow_timescale} summarizes the three time scales. The cloud horizon is selected to cover slow traffic interaction and lane-change evolution, typically 10--30 s. The vehicle-side emergency planner uses a shorter horizon, typically 5--10 s, to respond to abrupt local hazards and to keep a feasible escape branch when the most dangerous local scenario occurs. The controller uses a horizon of approximately 2--3 s, which covers the dominant liquid-sloshing and roll dynamics. This separation follows the service-time analysis in which cloud-side loops are affected by communication delay and tail latency, while the vehicle-side loop must satisfy emergency response and stability constraints \cite{xionglu2024_CVIS_survay,xuqing2021unreliable,wang2021platoonmpc}.

\subsection{Tank-Truck Rollover Risk Description}

The rollover risk of a tanker truck is driven by road curvature, lateral acceleration, articulation, roll motion, and liquid-sloshing forces. For a generic axle group, the load transfer ratio can be written as
\begin{equation}
  \mathrm{LTR}=\frac{F_z^R-F_z^L}{F_z^R+F_z^L},
  \label{eq:ltr_definition}
\end{equation}
where $F_z^R$ and $F_z^L$ are the right- and left-side vertical tire loads. The condition $|\mathrm{LTR}|=1$ indicates wheel lift-off. Because direct vertical tire-load measurement is difficult in engineering deployment, the controller uses a suspension-force-equivalent rollover index
\begin{equation}
  \mathrm{LTR}_{eql}
  =c_r(-k_{r1}\phi_1-c_1\dot{\phi}_1-k_{r2}\phi_2-c_2\dot{\phi}_2),
  \label{eq:ltr_eql_problem}
\end{equation}
where $\phi_i$ and $\dot{\phi}_i$ denote tractor or trailer roll angle and roll rate, $k_{ri}$ and $c_i$ are suspension roll stiffness and damping, and $c_r$ is a normalization coefficient. In the planning layer, simplified rollover models are used to screen candidate trajectories. In the control layer, $\mathrm{LTR}_{eql}$ is directly included as a constrained MPC output.

The selected trajectory must also account for liquid sloshing. Let $x_l$ denote the equivalent liquid state, including sloshing displacement or pendulum angles and their rates. The coupled tanker-truck state is expressed compactly as
\begin{equation}
  x=[x_v^T,\ x_l^T]^T,
\end{equation}
where $x_v$ contains longitudinal velocity, lateral velocity, yaw states, roll states, articulation-related states, and global pose. The coupled dynamics used by the vehicle-side planner and controller can be written as
\begin{equation}
  x_{k+1}=f(x_k,u_k,w_k),
  \label{eq:coupled_compact}
\end{equation}
where $w_k$ represents road, traffic, and modeling disturbances. Section~IV gives the control-oriented realization of this model.

\subsection{Hierarchical Planning and Control Problem}

The complete vehicle-cloud collaborative problem is formulated as a hierarchical optimization and arbitration problem:
\begin{equation}
  \tau_c^\star,d_c^\star
  =\arg\max_{\tau_c,d_c} \mathbb{E}\left[
  R_c(\mathcal{G}_{0:H_c},d_c,\tau_c)
  \right],
  \label{eq:cloud_problem}
\end{equation}
\begin{equation}
  \tau_v^\star
  =\arg\min_{\tau_v}
  J_v(\tau_v,\tau_c,\mathcal{P}_v)
  \quad
  \mathrm{s.t.}\quad
  g_v(\tau_v,\mathcal{P}_v)\leq 0,
  \label{eq:vehicle_planning_problem}
\end{equation}
\begin{equation}
  U^\star
  =\arg\min_{U}
  J_m(x_0,U,\tau_r)
  \quad
  \mathrm{s.t.}\quad
  g_m(x_0,U,\tau_r)\leq 0.
  \label{eq:mpc_problem_compact}
\end{equation}
In \eqref{eq:cloud_problem}, $R_c$ rewards safety, efficiency, comfort, rule compliance, and route consistency over the cloud prediction horizon $H_c$. In \eqref{eq:vehicle_planning_problem}, $J_v$ penalizes deviation from the accepted cloud guide, obstacle proximity, road-boundary violation, excessive curvature, and predicted rollover risk. In \eqref{eq:mpc_problem_compact}, $J_m$ penalizes tracking error, sloshing motion, rollover-index violation, actuator effort, and control increments.

The corresponding constraints include
\begin{equation}
  |\mathrm{LTR}_{eql,k}| \leq \mathrm{LTR}_{max}+\epsilon_{ltr,k},
  \label{eq:ltr_constraint_problem}
\end{equation}
\begin{equation}
  u_{min}\leq u_k\leq u_{max},\quad
  \Delta u_{min}\leq \Delta u_k\leq \Delta u_{max},
  \label{eq:actuator_constraint_problem}
\end{equation}
\begin{equation}
  r_k\in\mathcal{B}_{road},\quad
  \mathcal{B}_{veh}(r_k)\cap\mathcal{B}_{obs,k}=\emptyset,
  \label{eq:road_obstacle_problem}
\end{equation}
where $\epsilon_{ltr,k}$ is a slack variable for soft rollover constraints, $\mathcal{B}_{road}$ is the drivable road region, $\mathcal{B}_{veh}$ is the swept vehicle boundary, and $\mathcal{B}_{obs,k}$ is the obstacle occupancy set. This hierarchy does not require the cloud to control the vehicle directly. Instead, the cloud shifts risk avoidance earlier in time, the vehicle-side planner preserves a locally feasible fallback trajectory, and the MPC forms the final anti-rollover safety layer even when the selected reference is aggressive or the vehicle is already close to the rollover boundary.
